% !TeX root = ../navier-stokes-manuscript.tex
\subsection{The critical-height table}\label{sub:TheViscosityDrivenSpatialAscent}

The temporal Tower came from asking how a velocity can escape before a finite time.
The same index can now be read in the opposite direction.
Each time derivative of the critical height exposes a higher spatial Sobolev energy through viscosity.

For \(s\geq0\), set
\[
  E_s(t):=\frac12\norm{\Lambda^su(t)}_2^2
\]
and retain the complete nonlinear and pressure transfer in
\[
  \mathcal P_s(t)
  :=-
  \operatorname{Re}
  \pair{\Lambda^su}
  {\Lambda^s\bigl((u\cdot\nabla)u+\nabla p\bigr)}_{L^2}.
\]
The pressure contribution vanishes in this global scalar pairing because \(\Lambda^su\) is divergence-free, but keeping it in \(\mathcal P_s\) preserves the pressure-complete origin of the row.
Pairing the equation with \(\Lambda^{2s}u\) gives
\begin{equation}\label{eq:BodyFullSpatialEnergyBalance}
  E_s'=\mathcal P_s-2\nu E_{s+1}.
\end{equation}
One time derivative reaches one full Sobolev step higher.

At the critical height \(H=E_{1/2}\), the first row is
\begin{equation}\label{eq:BodyFirstCriticalRow}
  H'=\mathcal P_{1/2}-2\nu E_{3/2}.
\end{equation}
Across one Zeno first-passage interval, this becomes
\begin{equation}\label{eq:EachZenoPassageCost}
  \int_{t_k}^{t_{k+1}}\mathcal P_{1/2}(t)\,dt
  =L_k+2\nu\int_{t_k}^{t_{k+1}}E_{3/2}(t)\,dt.
\end{equation}
The nonlinear transfer has to supply both the next critical-height gain and the viscous cost accumulated during the same shrinking interval.
The identity does not yet prevent the transfer from becoming more intense as those intervals shorten.

Differentiate \eqref{eq:BodyFirstCriticalRow} and use \eqref{eq:BodyFullSpatialEnergyBalance} at the next height:
\begin{equation}\label{eq:BodySecondCriticalRow}
  H''
  =\partial_t\mathcal P_{1/2}
  -2\nu\mathcal P_{3/2}
  +(2\nu)^2E_{5/2}.
\end{equation}
One more derivative gives
\begin{equation}\label{eq:BodyThirdCriticalRow}
  H'''
  =\partial_t^2\mathcal P_{1/2}
  -2\nu\partial_t\mathcal P_{3/2}
  +(2\nu)^2\mathcal P_{5/2}
  -(2\nu)^3E_{7/2}.
\end{equation}
The pattern is triangular: every new derivative retains all lower transfer terms and opens one new highest spatial energy.

For \(n\geq1\), induction gives
\begin{equation}\label{eq:BodyCriticalHeightSpatialAscent}
  H^{(n)}
  =(-2\nu)^nE_{n+1/2}
  +\sum_{j=0}^{n-1}
  (-2\nu)^j
  \partial_t^{\,n-1-j}\mathcal P_{j+1/2}.
\end{equation}
The identity that isolates the top viscous term is
\begin{equation}\label{eq:BodyCompletedCriticalRow}
  \boxed{
  E_{n+1/2}
  =
  \frac{
    H^{(n)}
    -\displaystyle\sum_{j=0}^{n-1}
      (-2\nu)^j
      \partial_t^{\,n-1-j}\mathcal P_{j+1/2}
  }{(-2\nu)^n}}
  \qquad(n\geq1).
\end{equation}
Nothing has been discarded from the differentiated equation.
The higher spatial energy appears only after the full transfer history below it has been retained.

\begin{center}
\captionsetup{hypcap=false}
\captionof{table}{The temporal derivative of critical height and the higher spatial energy isolated in the same completed row.}
\label{tab:CompletedCriticalHeightRows}
\small
\renewcommand{\arraystretch}{1.45}
\begin{tabularx}{\textwidth}{@{}c >{\raggedright\arraybackslash}X c@{}}
\toprule
order & critical-height row & highest spatial energy \\
\midrule
\(0\) & \(H=E_{1/2}\) & \(E_{1/2}\) \\
\(1\) & \(H'=\mathcal P_{1/2}-2\nu E_{3/2}\) & \(E_{3/2}\) \\
\(2\) & \(H''=\partial_t\mathcal P_{1/2}-2\nu\mathcal P_{3/2}+(2\nu)^2E_{5/2}\) & \(E_{5/2}\) \\
\(3\) & \(H'''=\partial_t^2\mathcal P_{1/2}-2\nu\partial_t\mathcal P_{3/2}+(2\nu)^2\mathcal P_{5/2}-(2\nu)^3E_{7/2}\) & \(E_{7/2}\) \\
\(n\) & equation \eqref{eq:BodyCriticalHeightSpatialAscent} & \(E_{n+1/2}\) \\
\bottomrule
\end{tabularx}
\end{center}

The scalar derivative \(H^{(n)}\) is not the velocity row \(V_n\).
Its exact relation to the temporal Tower is the quadratic contraction
\begin{equation}\label{eq:BodyHeightDerivativeVelocityContraction}
  H^{(n)}
  =\frac12\sum_{a=0}^{n}\binom{n}{a}
  \operatorname{Re}
  \pair{\Lambda^{1/2}V_a}{\Lambda^{1/2}V_{n-a}}_{L^2}.
\end{equation}
Finite-time escalation reads the index \(n\) as another temporal response that cannot stay bounded.
The completed identity reads the same index as access to \(n+\tfrac12\) spatial derivatives of the velocity.
A concentrating singularity points toward rougher spatial behaviour; the viscous column points toward stronger Sobolev regularity.

This is the reversal at the centre of the argument.
The identities show its exact shape, but appearance is not control.
The next step is to show that the same spatial ascent exists inside every temporal row and then identify the finite object that can be bounded through \(T_*\).
